% !TEX root = ../monograph.tex
\chapter[Operator products and contact extensions]{Operator products and extensions across diagonals}
\label{ch18:products}

There are two distinct questions about colliding insertions. An operator
product expansion replaces a small cluster by finitely many local fields
when tested against separated spectators. A distributional extension
assigns values to test functions whose supports meet collision diagonals. The first requires
quantitative information about the marked renormalization group. The
second already follows, with a conservative finite order, from the
polynomial separated bounds. Neither construction defines a product of
point fields by simply setting their positions equal.

We use the common filtration $F_d$, its inclusions $\iota_{d,e}$,
and the mixed correlations $S$ from Chapter~\ref{ch13:continuum}.
The orthogonal actions on these spaces and the phase correlations
$S^\xi$ are those of Chapters~\ref{ch14:ward} and
\ref{ch15:phases}. Write $\mathfrak A$ for the
commutative compression algebra of Chapter~\ref{ch15:phases} and
$M(w)$ for the compression of a whole separated word $w$.
We first convert the lattice replacement estimate into the norm
bound needed for the operator product expansion. The contact extension constructed later does not use it.

The OPE originates with Wilson~\cite{Wilson1969}.
Fr\"ob and Holland establish the Yang--Mills OPE to every fixed order
of perturbation theory~\cite{FrobHolland2016}. That result explains the
relevant structure, but does not supply the nonperturbative estimate
required below. The standard distribution-extension argument is the
normal Taylor subtraction of Brunetti and Fredenhagen
\cite[Section~5.2, Theorems~5.2--5.3]{BrunettiFredenhagen2000}.
We use that procedure with an explicit choice adapted to the union of
all pair diagonals; no reconstruction theorem requires this particular
extension.

\section{A fixed prescription for all finite degrees}
\label{ch18:prescription}

A probe is a fixed word $v$ at distinct nonzero positions. It defines
$\lambda(A)=S(A(0)v)$ on each field space on which $A$ is defined.
These functionals separate $F_d$, by the definition of correlation-null
fields. Since $F_d$ is finite dimensional, finitely many probes span
its dual. A probe chosen at one degree continues to act at every higher
degree by the same correlation, using the physical inclusions.

One can choose fields $u_\alpha$ and finite linear combinations
$b_\alpha$ of probes, with a degree $|\alpha|$ attached to each, such that
\begin{equation}
 b_\alpha(u_\beta)=\delta_{\alpha\beta},\qquad
 A=\sum_{|\alpha|\le d}b_\alpha(A)u_\alpha\quad(A\in F_d),
 \label{ch18:dual}
\end{equation}
where fields in lower degrees are included in $F_d$ when necessary.
All these identities hold simultaneously at every finite degree.
Here is the useful linear-algebra construction. Suppose the choices
through $d$ have been made and let $R:F_d\to F_{d+1}$ be inclusion.
The old probes give a map $B:F_{d+1}\to F_d$ with $BR=1$.
For any section $S_0$ of the quotient map
$q:F_{d+1}\to F_{d+1}/RF_d$, replace it by
$S_0-RBS_0$. The corrected section is killed by every old probe.
Lift a quotient basis through it, and represent its dual quotient
coordinates by finite probe combinations. This extends
\eqref{ch18:dual} without changing an earlier field or probe.
In particular $b_\alpha$ vanishes on all degrees below $|\alpha|$.
This is only finite-dimensional linear algebra; no infinite expansion
of a field is being asserted.

For a cluster word $w$, define the unsymmetrized field reconstruction
\[
 H_d(w)=\sum_{|\alpha|\le d}b_\alpha(w)u_\alpha,
\]
where $b_\alpha(w)$ means replacing the distinguished field in its
probe correlations by the entire cluster. The mixed table is always
used for this measurement. Thus the coefficients have no dependence
on a phase or on spectators later placed outside the cluster.
At a center $o$, apply this definition to $\tau_{-o}w$.
For each finite collection of probes there is a positive clearance:
a cluster sufficiently close to $o$ cannot meet a probe point, even
after rotating the cluster. All subsequent formulas are used on this
clear region.

Write $R_d(K)$ for the orthogonal action on $F_d$, with normalized
Haar measure on $O(4)$. Symmetrize the reconstruction by
\[
 A_dw=\int_{O(4)}R_d(K)H_d(K^{-1}w)\,dK,
 \qquad B_dw=\tfrac12\{A_dw+\overline{A_d\bar w}\}.
\]
Both maps fix every single field of degree at most $d$ inserted at zero.
They are smooth on the clear configuration region, since finitely many
smooth correlation functions are integrated over a compact group.
They preserve reality and intertwine rotations. Translation covariance
is implemented by translating the center as above.

Averaging must be followed by a compatible removal of lower degrees.
Let $I_qA=A(0)$, put $L_{i,q}=B_iI_q$, and define
\begin{equation}
 Q_0=1,\qquad
 Q_j=1-\sum_{i<j}\iota_{i,j}Q_iL_{i,j},\qquad
 E_j=\bigcap_{i<j}\ker(Q_iL_{i,j}).
 \label{ch18:graded}
\end{equation}
The inclusion identities give, by induction,
\[
 Q_jL_{j,q}=0\quad(q<j),\qquad
 \sum_{i\le j}\iota_{i,j}Q_iL_{i,q}=\iota_{q,j}\quad(q\le j).
\]
Consequently $Q_j$ is the projection onto $E_j$ along
$\iota_{j-1,j}F_{j-1}$. These spaces are invariant under rotations
and conjugation. Fix real bases $(O_{j,a})_a$ of $E_j$, taking the
unit itself in degree zero, and denote their coordinate functionals
by $\omega_{j,a}$. The coefficients used in this chapter are
\begin{equation}
 c_{j,a}(w;o)=\omega_{j,a}Q_jB_j(\tau_{-o}w).
 \label{ch18:coefficients}
\end{equation}
The fields, probes and projections are fixed once, before any
truncation, phase, spectator list or small-distance limit.


\section{The marked replacement estimate}
\label{ch18:input}

The expansion of Chapter~\ref{ch11:insertions} keeps every joint
nilpotent marker coefficient. To replace a cluster by a single field,
one needs more than the one-insertion graph contraction. Markers may
merge as their separation becomes smaller than a block. Their source
weights multiply; an incoming term involving at least two insertions is kept as one factor,
with its support, prescribed set of large-field sites and eligibility condition. Its stronger support
norm absorbs the loss when several marks are assigned to the same coarse
region. After local Taylor subtraction, the signed residual must be
transported on the same history and compared in the physical quotient.
Taking absolute values before that subtraction loses the cancellation.

Here is the quantitative replacement estimate. It is formulated
so that its two uses, determination of local polynomial coefficients
and phase-uniform error control, are visible. Put $w_0=\sum_i d_i$ for a fixed cluster degree
list, and let
\[
 x_i=o+r z_i,\qquad |z_i|\le M,\qquad |z_i-z_j|\ge\eta>0.
\]
Cluster labels range over fixed bounded subsets of their finite-dimensional
spaces. Spectator labels have bounded degrees and coefficients; their
positions stay at distance at least $c>0$ from $o$ and from one another.
The constants below may depend on these finite lists, $M,\eta,c$ and
a requested error exponent, but not on positions, phases or $r$.

\begin{proposition}[Same-history marked replacement]
\label{ch18:replacement-input}
For every finite degree $D$ and every $\epsilon>0$, on the geometry
just specified and for all sufficiently small $r$, the unsymmetrized reconstruction
coefficients obey
\[
 |b_\alpha(w;o)|\le C r^{|\alpha|-w_0-\epsilon}.
\]
For each spectator word $v$, set
\[
 Z_D=M(wv)-M((H_D(\tau_{-o}w))(o)\,v).
\]
The actual lattice marked replacement, with its signed residual and
all original terms involving at least two insertions retained, gives the bound uniform in the number of reflected copies
\begin{equation}
 \langle\Omega,(Z_D^*Z_D)^k\Omega\rangle
       \le (C r^{D+1-w_0-\epsilon})^{2k},\qquad k\ge1.
 \label{ch18:replacement-moments}
\end{equation}
The constant and the small-$r$ threshold precede $k$; cutoff thresholds
may depend on each fixed $k$. The same estimates hold for every fixed
finite list of differentiated insertions, with the corresponding
scaling dimensions, and for the finitely many probe families above,
uniformly under their rotations and conjugations.
\end{proposition}

This is an analytic estimate on the selected family, not an assertion
that any bounded marked expansion has an OPE. Its construction requires
the signed local transport, the stronger anchored marked support bound,
and the graded quotient comparison. In particular, if a moving field
has coefficient size $\varepsilon_m$ at scale $m$, the needed quotient
estimate is of the form
\[
 \operatorname{dist}_{F_t}(A,\iota_{i-1,t}F_{i-1})
       \le a_m^{-(t-i)}e^{o(\log(1/a_m))}\varepsilon_m .
\]
The lower field must be identified by all spectator responses, using the
normalized-null lemma of Chapter~\ref{ch13:continuum}. A bound on a chosen
coordinate representative alone does not supply this statement.
The estimate uniform in the number of reflected copies also cannot be replaced by a bound on the mixed
mean of $Z_D$. These are supplied by the lattice estimate of Section~\ref{ch18:lattice-products} and its passage to the compression algebra below.

% !TEX root = ../monograph.tex
\begin{proof}[Proof of Proposition~\ref{ch18:replacement-input}]
\label{ch18:raw-replacement}
The lattice estimate in Proposition~\ref{ch18:lattice-replacement}
was proved before phase decomposition. We use its uniformity in
replica number to obtain the algebra-valued statement for the
reconstructions $H_D$ fixed above.

Let $F$ be an interior lattice difference with chessboard seminorm
at most $B$, and let its exterior factors have interior seminorms
at most $B_1,\ldots,B_q$. Put copies of the complete difference
and exterior factors in separated cells before taking limits.
The repeated-word identity and chessboard inequality give, for its
compressed limit $Z$ and each integer $n\ge1$,
\[
 \int_\Xi|\xi(Z)|^{2n}\,d\nu(\xi)
                   \le(B\prod_{l=1}^qB_l)^{2n}.
\]
The seminorm bounds $B,B_1,\ldots,B_q$ are independent of $n$.
For each fixed $n$, the cutoff and the separation of the copies may
need to be increased.
This uses the complete signed density, with every term involving at least two insertions,
from the proof of \eqref{ch18:lattice-estimate}.
Equation~\eqref{ch15:moment-norm} then gives
$\|Z\|\le B\prod_lB_l$. Positivity here belongs to the limiting
commutative algebra and the original Wilson law, not to an
individual signed effective activity.

Apply this observation to the auxiliary field $A_t$ in the proof
of Proposition~\ref{ch18:mixed-bound}. On grid configurations its
replacement error has norm with exponent
$Z_0-(t+1/16)(\theta-\theta_-)$, where
$Z_0=W\theta_++4n_{\rm cl}(\theta_+-\theta)$.
Fixed exterior factors cost a bounded factor. The graded quotient
bound gives exponent $Z_0-i\theta$ for a degree-$i$ component of
$A_t$; all the finitely many fields used in its reconstruction
have uniformly bounded compressed exterior words. Testing the
same error against the fixed probes gives the probe coordinate
errors. With $t$ chosen as in that proof, the omitted degrees
and the low-coordinate errors therefore give
\[
 \|M(wv)-M((H_D(\tau_{-o}w))(o)v)\|
                \preccurlyeq e^{[Z_0-(D+1)\theta]T}.
\]
There is no spectator-dependent choice of $H_D$: the only
geometry-dependent auxiliary field has disappeared from this
inequality.

The derivative descent of Proposition~\ref{ch18:mixed-bound}
works in this norm. Equivalently, apply it to every character,
using the same sharp chessboard seminorm bounds, and take their supremum.
The crude and grid constants are uniform in the character.
A positional derivative changes its input degree by one; a finite
descending induction removes the grid restriction, also for any
fixed list of differentiated fields. Letting the two resolution
heights tend to one yields
\[
 \|M(wv)-M((H_D(\tau_{-o}w))(o)v)\|
                      \le C r^{D+1-W-\varepsilon}.
\]
The coefficient assertion is exactly the corresponding fixed-probe
estimate in Proposition~\ref{ch18:mixed-bound}. Compact label
sets, conjugation, and the finite probe families do not change
its constants. For their rotation orbits the same uniform estimates
apply on the compact set of rotated labels and positions; choose
$r$ so all these probes remain separated from the cluster.

Finally apply the vacuum state to $(Z_D^*Z_D)^n$. Its norm is
at most $\|Z_D\|^{2n}$, giving
\eqref{ch18:replacement-moments} with the same $C$ and the same
small-$r$ threshold for every $n$. Thus the order of quantifiers
is obtained from the norm estimate, rather than from an
uncontrolled sequence of separate moment constants.
\end{proof}


\section{Finite-order operator product expansion}
\label{ch18:ope}

\begin{theorem}[Phase-uniform finite-order expansion]
\label{ch18:ope-theorem}
The fixed coefficients
\eqref{ch18:coefficients} satisfy, for every finite $D$ and
$\epsilon>0$,
\begin{align}
 |c_{j,a}(w;o)|&\le C r^{j-w_0-\epsilon},\qquad j\le D,\notag\\
 \left\|M(wv)-\sum_{j\le D,a}c_{j,a}(w;o)
                        M(O_{j,a}(o)v)\right\|
      &\le C r^{D+1-w_0-\epsilon}.
 \label{ch18:ope-bound}
\end{align}
The constants and threshold are uniform on the bounded label sets and
separated geometries of Section~\ref{ch18:input}. Evaluating a character,
or the mixed vacuum functional, gives the same bound for its correlations.
\end{theorem}
\begin{proof}
Full support of the phase measure converts
\eqref{ch18:replacement-moments} into the operator-norm estimate by
\eqref{ch15:moment-norm}. To compare the unsymmetrized and symmetric reconstructions,
write $h=H_H(\tau_{-o}w)$ at an auxiliary finite degree $H$.
Every component of $B_jw-L_{j,H}h$ is an average of finitely many
replacement errors tested against rotated probe words. Their clearance
is uniform by compactness of $O(4)$. Hence
\[
 |B_j(\tau_{-o}w)-L_{j,H}h|
                         \le C r^{H+1-w_0-\epsilon}.
\]
This identity uses covariance of the mixed correlations to rotate the
probes, as well as the cluster; rotating labels alone would not suffice.

Apply $Q_j$. In the expansion of $h$, every field of degree below $j$
is killed. The remaining finitely many terms have powers at least
$r^{j-w_0-\epsilon}$, by the raw coefficient estimate. Taking $H\ge j$
proves the first bound. For the second, take $H=D$ and sum the reconstruction
comparison through $j=D$. The reconstruction identity following
\eqref{ch18:graded} returns $h$ exactly. Multiplying the coefficient
errors by $M(O_{j,a}(o)v)$ costs a uniform constant on the separated
spectator geometry. Add the raw operator error. Finitely many sums and
averages are controlled by first using a smaller error exponent.
\end{proof}

For two fields of degrees $d_A,d_B$, the degree-$j$ coefficient thus
has size at most $C r^{j-d_A-d_B-\epsilon}$. The remainder may still
grow when $D+1<d_A+d_B$; its order nevertheless improves with $D$.
No assertion concerns the convergence of the infinite sum in $D$.

The derivative version uses the differentiated marked replacement estimate. One first fixes the finite derivative order and the degree
margin, then carries out the same finite probe functionals. The derivative
identities of the field filtration identify the resulting derivatives.
Uniform differentiated bounds permit passage from a dense grid to all
separated positions by the mean-value theorem. Mere pointwise grid
estimates without those bounds would not justify that passage.

For two nested clusters at radii $\rho<R$, replacements can be composed
by telescoping the signed difference. The first error is tested against
the unchanged second cluster; the next is multiplied by the already
produced coefficients. Accordingly a uniform nested estimate requires
the replacement bound when exterior separation also tends to zero,
with its dependence on $R$ retained. Fixed-positive-clearance
Proposition~\ref{ch18:replacement-input} by itself does not give that
stronger result. For any prescribed finite tree, the required additional
statement is a signed bound at every vertex, with degree-$q$ coefficient
cost $\rho^{q-d}$ and residual gain $(\rho/R)^{D+1}$, up to arbitrary
small power losses, including the coefficients and the operators that its descendants apply to the test function. With those bounds, finite telescoping proves the
corresponding finite-tree expansion. This stronger multi-scale assertion
is not needed for the distributional extension below.

\section{An explicit finite-order extension across every diagonal}
\label{ch18:extension}

We now use only a uniform polynomial separated bound. Fix an arity
$n\ge2$ and a field list $\boldsymbol A$. Write
\[
 \Delta=\bigcup_{i<j}\{x_i=x_j\},\qquad
 \delta(x)=\min_{i<j}|x_i-x_j|.
\]
Suppose, as supplied by the separated bounds and the uniform phase slab
bounds, that for some finite $P\ge0$,
\begin{equation}
 |S_n^\xi(\boldsymbol A;x)|\le C_K\min(1,\delta(x))^{-P}
 \quad(x\in K\setminus\Delta),
 \label{ch18:polynomial}
\end{equation}
uniformly in $\xi$ on every compact $K$. The mixed family has the same
bound. No optimal value of $P$ is required. Choose once an integer
$k\ge0$ with $k+1>P$, and a length $s>0$.

We specify which normal jet is subtracted near simultaneous collisions.
Let $m=n(n-1)/2$ be the number of pairs and put, on $\Conf_n$,
\[
 \rho(x)=\left(\sum_{i<j}|x_i-x_j|^{-2}\right)^{-1/2}.
\]
Then $\delta/\sqrt m\le\rho\le\delta$.
Choose a fixed nonnegative smooth function $\varphi$ on $[0,\infty)$,
equal to one on $[0,2\sqrt m]$ and zero on $[4\sqrt m,\infty)$.
For a pair $h=\{i,j\}$ define
\begin{equation}
 a_h(x)=\frac{\varphi(|x_i-x_j|/\rho(x))}
                   {\sum_{u<v}\varphi(|x_u-x_v|/\rho(x))}.
 \label{ch18:pair-weights}
\end{equation}
The denominator is at least one, because a closest pair has ratio at
most $\sqrt m$. Thus $a_h\ge0$ and $\sum_h a_h=1$.
Only pairs whose separation is at most $4\sqrt m\rho$ are selected.
These weights are smooth off $\Delta$, invariant under simultaneous
translations and rotations, and permuted with the slots.

Let $\pi_hx$ replace $x_i,x_j$ by their centroid, leaving the other
positions unchanged, and put $z_h=x-\pi_hx$. This is the orthogonal
normal displacement from the pair diagonal; its norm is
$|x_i-x_j|/\sqrt2$. Define the normal Taylor polynomial of a test function by
\begin{equation}
 T_{h,k}f(x)=\sum_{\ell=0}^k\frac1{\ell!}
             D^\ell f(\pi_hx)[z_h,\ldots,z_h].
 \label{ch18:normal-Taylor}
\end{equation}
Only normal directions occur; the centroid and all external positions
remain parameters. Finally fix $\chi\in C^\infty([0,\infty))$ with
$0\le\chi\le1$, $\chi=1$ on $[0,1]$ and $\chi=0$ on $[2,\infty)$,
and set
\begin{equation}
 (\mathcal W f)(x)=f(x)-\chi(\rho(x)/s)
                                  \sum_h a_h(x)T_{h,k}f(x).
 \label{ch18:subtraction}
\end{equation}
This formula, including $k,s,\varphi,\chi$, is the subtraction convention.
It is the same in every phase and is fixed before the test function.
It deliberately subtracts more jets than an optimal scaling-degree
prescription may require.

\begin{theorem}[A common contact extension]
\label{ch18:extension-theorem}
Under \eqref{ch18:polynomial}, the formula
\begin{equation}
 \mathfrak D_{\boldsymbol A}^\xi(f)
   =\int_{\Conf_n} S_n^\xi(\boldsymbol A;x)
                                      (\mathcal Wf)(x)\,dx
 \label{ch18:extension-formula}
\end{equation}
defines a distribution of local order at most $k+1$.
It extends the separated correlation, is uniform in the phase on
compact supports of test functions, and is continuous in $\xi$ on each test function.
Averaging over phases gives the same formula for the mixed correlations.
Translations, orthogonal covariance and simultaneous permutation
symmetry are preserved by this prescription.
\end{theorem}
\begin{proof}
This is the normal Taylor extension argument of
\cite[Theorem~5.3]{BrunettiFredenhagen2000}; the following verification
explains why one subtraction treats every partial diagonal here.
When $\rho\le s$, the weights sum to one, so
\[
 \mathcal Wf=\sum_h a_h(f-T_{h,k}f).
\]
For each selected pair, $|z_h|\le 2\sqrt{2m}\rho$.
Taylor's integral remainder therefore bounds this expression by
$C\rho^{k+1}\|f\|_{C^{k+1}}$. Since $\delta\ge\rho$,
\eqref{ch18:polynomial} bounds its product with the kernel by
$C\rho^{k+1-P}\|f\|_{C^{k+1}}$, which is bounded and hence locally
integrable. This remains true when arbitrarily many other separations
vanish at independent rates: the selected pair always has displacement
comparable to the smallest separation. No fixed collision pattern or
single radial approach is assumed.

On the transition region $s\le\rho\le2s$, all separations are bounded
below by $s$, and the finite Taylor polynomials are bounded by the same
seminorm of the test function. Outside this region the integrand uses $f$ itself. If $f$ is supported
in a compact $K$, a nonzero subtracted jet requires $\pi_hx\in K$ and
$|z_h|\le C_ns$. Thus the integrand is supported in one fixed compact
enlargement $K'$ of $K$. We obtain
\[
 |\mathfrak D_{\boldsymbol A}^\xi(f)|
        \le C_{K'}\max_{|\alpha|\le k+1}
                              \|\partial^\alpha f\|_\infty.
\]
This proves continuity and the local order assertion. If $f$ is supported
away from $\Delta$, all its jets at every $\pi_hx\in\Delta$ vanish,
so \eqref{ch18:extension-formula} is the original separated integral.
Continuity in the phase and the mixture formula follow by dominated
convergence with the same bound. The normal projections, pair weights
and radial cutoff commute with the stated geometric symmetries.
\end{proof}

For $n=1$ use the original smooth one-point function, and for $n=0$
use the normalized unit. A consistent multilinear family on the whole
filtration is obtained using the fixed invariant summands $E_j$ in
\eqref{ch18:graded}. For each finite tuple $(j_1,\ldots,j_n)$ choose
$P$ uniformly on the finite basis lists of $E_{j_1}\times\cdots\times
E_{j_n}$, and choose the same $P,k$ for every permutation of that tuple.
Use this one subtraction on the tuple and extend multilinearly there.
Every field has a unique finite sum of components in the $E_j$, so
extension by finite multilinearity gives a prescription on all of $F$.
It is compatible with the inclusions and preserves rotations because
each $E_j$ is invariant. Choosing an order independently for each vector,
or redefining it on each larger degree space, would not have these
properties.

\section{Sharper subtraction and dependence on the prescription}
\label{ch18:scope}

The explicit extension above is conservative. A degree-sensitive
alternative assigns to every collision subset $h$ the normal dimension
$q_h=4(|h|-1)$, the total input degree $d_h=\sum_{i\in h}d_i$, and
\[
 L_h=\max(0,d_h+1-q_h).
\]
If a contact projection extracts a degree-$q$ field from that cluster,
one subtracts the normal Taylor jet of the test function through order $L_h-1-q$ when
$q<L_h$. A single shell then has the combined power
\[
 \rho^{q-d_h}\,\rho^{L_h-q}\,\rho^{q_h}
                  =\rho^{q_h+L_h-d_h},
\]
with exponent at least one. This explains the source prescription and
its proposed sharper local-order bound. For several collisions the
shell estimate must be established for the complete signed expression,
including nested subtractions and all child operators acting on test functions. The parent
Taylor operator acts first on the test function; the child operators act after
it. They cannot be commuted for nested subsets. A partition of unity
cannot be moved through these normal Taylor operators without its
resulting diagonal jets. A one-shell estimate alone does not verify
that stronger forest prescription. We have therefore used the fully
specified pair-weighted subtraction \eqref{ch18:subtraction}, for which
absolute convergence at all collision patterns was proved directly.

Two extensions agreeing on $\Conf_n$ differ by a distribution supported
on $\Delta$. Along a smooth stratum, a finite-order difference is a sum
of normal derivatives of delta distributions with distributional
coefficients. This is the usual local renormalization freedom
\cite[Section~5.2]{BrunettiFredenhagen2000}. Changing the cutoff or
subtraction order can change contact terms while preserving separated
correlations and the OPE estimates. Contact reflection positivity,
contact Ward identities, uniqueness, unit removal and identification
with unsubtracted coincident-point lattice limits require additional
normalization arguments.
